\subsection*{MT-M27. Full adelic trace class on a wide strip}
Let $K$ be a fixed number field, $C=C_K$, and
$C^1=\ker|\cdot|$. Choose a norm splitting and normalize
the compact Haar measure to have volume one. Put
\[
 \begin{split}
 E_K&=\mathcal S_{\mathrm E}(C),\qquad
 \Theta_K=\Phi_K(\mathcal S_0(\mathbf A_K)),\\
 H_{J,s}&=L^2\!\left(C,
       \max\{|g|^{2c},|g|^{2d}\}
       (1+|\log|g||^2)^{s/2}\,dg\right),\\
 \mathcal J_{J,s}&=\overline{\Theta_K}^{\,H_{J,s}},
 \qquad Z_{J,s}=H_{J,s}/\mathcal J_{J,s},
 \end{split}                                                    \tag{M27.1}
\]
where $s\ge0$ and $J=[c,d]$ satisfies $c<0<1<d$.
The topology on $E_K$ is the one specified in Masullo
Definitions D.3.3.1/D.3.3.3, original PDF pp.~804--805.
The Hilbert space and denominator in (M27.1) are his
specified norm and Hilbert closure, before the harmless
centering twist of MT-M24.

For each $\chi\in\widehat{C^1}$ choose a unitary extension
$\beta_\chi$ to $C$ and use the Hecke convention of MT-M22.
Let $\Xi_\chi$ have zeros $\rho$ with multiplicities
$m_{\chi,\rho}$. Define
\[
 \widehat f(\beta_\chi|\cdot|^z)
       =\int_C f(g)\beta_\chi(g)|g|^z\,dg.
                                                               \tag{M27.2}
\]
Then every $f\in\mathcal D(C)$ induces a trace-class
convolution operator on $Z_{J,s}$, and
\[
 \operatorname{Tr}_{Z_{J,s}}C_f
     =\sum_{\chi\in\widehat{C^1}}
        \sum_{\Xi_\chi(\rho)=0}
                   m_{\chi,\rho}
                     \widehat f(\beta_\chi|\cdot|^\rho).
                                                               \tag{M27.3}
\]
The series is absolutely convergent, uniformly on bounded
subsets of $\mathcal D(C)$. The trace norms are uniformly
bounded on those same sets. This includes the infinitely
many archimedean characters at any fixed finite level.
No GRH is required.

There are also exact kernel identities:
\[
 \begin{split}
 \overline{\Theta_K}^{\,E_K}
   &=\{u\in E_K:
      \partial_z^k\widehat u(\beta_\chi|\cdot|^z)|_{z=\rho}=0
            \text{ for all }\chi,\rho,0\le k<m_{\chi,\rho}\},\\
 \mathcal J_{J,s}
   &=\{u\in H_{J,s}:
      \partial_z^k\widehat u(\beta_\chi|\cdot|^z)|_{z=\rho}=0
            \text{ for all }\chi,\rho,0\le k<m_{\chi,\rho}\}.
 \end{split}                                                    \tag{M27.4}
\]
Thus the corrected full smooth quotient
$E_K/\overline{\Theta_K}^{E_K}$ maps continuously and
injectively, with proper dense image, into the specified
Hilbert quotient. We do not identify its stronger smooth
topology with the Hilbert norm topology.

\paragraph{Classical input and exact scope.}
The additional uniform input here is the published Meyer
Lemma~5.6, author PDF pp.~41--42, in
\emph{On a representation of the idele class group related
to primes and zeros of $L$-functions},
\url{https://arxiv.org/pdf/math/0311468} (v3).
For sufficiently large finite $S$, restricting its
stated $C_K$-valued map to $U_S$-invariants, as in
its proof, gives
\[
 (P_--P_+^S)C_f:
      L^2(C_S)_\alpha\oplus L^2(C_S)_{1-\alpha}
           \longrightarrow
          \mathcal S(C_S)_{>}\oplus\mathcal S(C_S)_{<}
                                                               \tag{M27.5}
\]
extends to a nuclear operator for $\alpha>1$,
uniformly for $f$ in bounded test sets.
Here $>$ means all weights $>1$, $<$ all weights $<0$,
$P_-$ is the diagonal projection, and $P_+^S$ has range
in the embedded adelic summation space $H_+^S$.
The supplied proof of this lemma was read in full,
including its Poisson and cutoff reduction; its cited
upstream local kernel lemmas are classical inputs, not
claimed newly fully audited. The rank-one meaning of
``nuclear'' is explicitly given in Meyer \S2.1, formula
(5), author pp.~6--7. We prove the transfer from (M27.5)
to the particular Hilbert quotient; it does not follow
merely from equality of spectral multiplicities.

\begin{proof}
\emph{Character projections and the two closed kernels.}
The compact abelian Fourier decomposition gives an
orthogonal Hilbert sum
\[
                     H_{J,s}=\mathop{\widehat\bigoplus}_\chi
                                      H_{J,s}(\chi).
                                                               \tag{M27.6}
\]
The weight depends only on the norm, so the compact
character projectors are orthogonal contractions.
The subspace $\Theta_K$ is invariant under these
projectors: by the global set identity (M24.7), it is
the intersection of the two embedded Meyer spaces,
both preserved by compact averaging. Its $\chi$ image
is exactly $\Theta_\chi$ of MT-M24.
Consequently $\mathcal J_{J,s}\cap H_{J,s}(\chi)$ is
the Hilbert closure of $\Theta_\chi$: apply the
contractive projector to a convergent approximating
sequence from $\Theta_K$. MT-M26 identifies this sector
closure with its full zero-jet kernel, because
$[0,1]\subset\operatorname{int}J$.
All these jet functionals are continuous on the full
Hilbert space: the compact projection has norm one,
and the Cauchy--Schwarz estimate (M26.7) applies in
that sector. Finite sums of character projections
converge strongly in (M27.6). It follows that vanishing
of all character/zero jets is equivalent to membership
in $\mathcal J_{J,s}$. This proves the second identity
in (M27.4).

For completeness the corresponding smooth closure
requires no uniform inverse estimates. An element
$u\in E_K$ lies at one elementary compact-invariance
level, with compact factor a finite abelian group
times a torus. At this level its seminorms are the
weighted suprema of all derivatives in the norm and
compact directions. Integration by parts on that
torus bounds every Fourier coefficient, in each
fixed norm-direction seminorm, by arbitrarily high
negative powers of its lattice frequency. For
example, apply $(1-\Delta)^N$ in the compact variable;
the character eigenvalue grows as $(1+|\nu|^2)^N$.
Taking $2N$ larger than the torus dimension makes
the Fourier tails summable, and increasing $N$
treats any specified compact derivative as well.
Thus finite character truncations converge to $u$
in every seminorm of that elementary level, hence
in $E_K$.
If all its zero jets vanish, each single-character
truncation belongs to the actual image $\Theta_\chi$
by MT-M24. Every finite truncation is therefore in
$\Theta_K$, and $u\in\overline{\Theta_K}^{E_K}$.
The reverse inclusion follows from continuity of
each zero-jet functional and its vanishing on
$\Theta_K$. This proves the first identity in
(M27.4). The natural map from the corrected smooth
quotient into the Hilbert quotient is consequently
injective and continuous. Its image is dense:
finite character sums times smooth compactly
supported functions of $\log|g|$ belong to $E_K$
and are dense in (M27.6). It is proper, since
surjectivity would, after a single compact-character
projection, contradict the proper-image assertion
of MT-M26 in that fixed sector.

\emph{A finite level for each test set, not a finite set of characters.}
Choose $\alpha>1$ so that
$[1-\alpha,\alpha]\subset(c,d)$.
A test $f\in\mathcal D(C)$ is invariant under
$U_S=\prod_{v\notin S}\mathcal O_v^\times$
for sufficiently large finite $S$ containing all
infinite places; $S$ can be enlarged to make the
$S$-ideal class group trivial. One common $S$ works
for a bounded set of tests. This follows directly
from the compact-invariance-level definition of
the test-function bornology, and is stated
explicitly in Meyer's proof of Theorem~5.8,
author p.~43. Enlarging $S$ preserves invariance.
For such $S$,
$C/U_S\simeq C_S=\mathbf A_S^\times/\mathcal O_{K,S}^\times$:
the trivial $S$-class group lets one adjust all
valuations outside $S$ by a principal idele; the
kernel is precisely the $S$-unit group.

Averaging over $U_S$ is an orthogonal contraction
on $H_{J,s}$, preserves $\mathcal J_{J,s}$, and
therefore gives an orthogonal projection on
$Z_{J,s}$. Since $f$ is $U_S$-invariant,
$C_f=C_fP_{U_S}=P_{U_S}C_f$.
It is enough to prove trace class on this invariant
Hilbert subspace, identified with functions on
$C_S$, and then extend it by zero on the orthogonal
complement. The invariant denominator is the
closure of $P_{U_S}\Theta_K$, by the same averaging
argument as above. No truncation of the infinitely
many compact Lie characters of $C_S^1$ is made.

\emph{An explicit bounded map from Meyer's ambient pair.}
Write
\[
 A_S=\mathcal S(C_S)_{>}\oplus\mathcal S(C_S)_{<},
 \qquad E_S=\mathcal S(C_S)_{\mathbf R},
 \qquad i_-(h)=(h,h).
\]
Choose a smooth norm cutoff $\phi$ equal to one
near zero and zero near infinity. Meyer's map
\[
        p_-(u_0,u_1)=(1-\phi)u_0+\phi u_1
                                                               \tag{M27.7}
\]
is bounded from $A_S$ to $E_S$.
This also follows immediately from weighted
seminorms: on the upper tail use a right weight
larger than the desired weight, on the lower
tail a smaller left weight, and treat derivatives
of the cutoff on a compact norm interval.
It satisfies $p_-i_-h=h$ and $P_-=i_-p_-$.

Put $\mathcal P=\operatorname{span}\{1,|\cdot|^{-1}\}$.
Poisson gives
\[
 X_S=H_+^S+i_-(E_S)
      =\{(u_0,u_1)\in A_S:u_0-u_1\in\mathcal P\}.
                                                               \tag{M27.8}
\]
Here is the explicit splitting needed for this identity
and for the subsequent bounded map. Choose two adelic
Schwartz functions, invariant outside $S$, whose pairs
$(f(0),\widehat f(0))$ are linearly independent.
They exist by taking fixed standard finite factors
and archimedean smooth functions: one vanishes at
zero but has nonzero integral, and another has
nonzero value at zero. Normalize their two pairs.
Their embedded images in $H_+^S$ give a linear map
$V:\mathcal P\to H_+^S$ such that
$dV=\operatorname{id}_{\mathcal P}$, where
$d(u_0,u_1)=u_0-u_1$. This uses exactly Poisson
(M24.6), including its two signs.
For any pair in $A_S$, evaluate its difference at
norms $1$ and $e$, and let $\ell(u_0,u_1)$ be the
unique element of $\mathcal P$ with those two
values. The matrix
$\left(\begin{smallmatrix}1&1\\1&e^{-1}\end{smallmatrix}\right)$
is invertible, so $\ell:A_S\to\mathcal P$ is a
bounded linear map and equals $d$ on $X_S$.
Subtracting $Vd$ makes the two entries equal;
their common function has all real Schwartz
weights. This proves (M27.8) directly.

Let $q:H_{J,s}^{U_S}\to Z_{J,s}^{U_S}$ be the
Hilbert quotient map. Inclusion $E_S\to H_{J,s}^{U_S}$
is bounded, by the exponential weights and compactness
of the norm-one factor. Therefore
\[
 B:A_S\longrightarrow Z_{J,s}^{U_S},\qquad
 B(y)=q\bigl(p_-(y-V\ell(y))\bigr)
                                                               \tag{M27.9}
\]
is bounded. It satisfies $B(i_-h)=q(h)$.
It vanishes on $H_+^S$: if $y\in H_+^S$, then
$y-V\ell(y)$ is still in $H_+^S$ and is diagonal;
its common function belongs to the theta image
by the Poisson intersection identity (M24.7),
so its Hilbert quotient class is zero.
Thus on $X_S$, (M27.9) is the actual canonical
quotient map, with an explicitly defined bounded
extension to the whole ambient pair. No assertion
about an unspecified quotient topology is used.

\emph{Transfer of the nuclear extension.}
There is a bounded diagonal inclusion
\[
 \delta:H_{J,s}^{U_S}\longrightarrow
       L^2(C_S)_\alpha\oplus L^2(C_S)_{1-\alpha},
       \qquad h\longmapsto(h,h).
                                                               \tag{M27.10}
\]
Indeed both endpoint exponential weights on the
right are bounded by the maximum of the weights
at $c,d$; the polynomial factor in (M27.1) is
at least one. The inclusion has norm at most
$\sqrt2$ for the chosen normalizations.
Compose (M27.5) with this inclusion and (M27.9):
\[
              A_f=B(P_--P_+^S)C_f\delta:
                     H_{J,s}^{U_S}\to Z_{J,s}^{U_S}.
                                                               \tag{M27.11}
\]
It is a nuclear map between Hilbert spaces,
uniformly for the bounded test sets under consideration.
To make the bornological-to-Hilbert step explicit,
Meyer's definition writes the middle map as
$\sum_j\lambda_j w_j\ell_j$, where
$\sum|\lambda_j|<\infty$, $(w_j)$ is bounded in
$A_S$ and $(\ell_j)$ is equibounded in the source
dual. The images $Bw_j$ are norm bounded.
The functionals $\ell_j\delta$ are uniformly
norm bounded. If a precompact rather than the
norm-bounded convention is used on the Hilbert
source, the same assertion follows from uniform
boundedness: equiboundedness on its compact
subsets implies pointwise boundedness of these
continuous functionals, hence a common operator
norm bound. A linear functional bounded on all
compact subsets of a Banach space is continuous,
since otherwise a sequence tending to zero
could have unbounded images. This also verifies
that no bornology convention inserts a hidden
topology assumption. The resulting rank-one
series has summable products of Hilbert norms.

On a smooth $h\in E_S$, the definitions give
$P_-C_fi_-h=i_-C_fh$, while
$P_+^SC_fi_-h\in H_+^S$ is killed by $B$.
Consequently $A_fh=q(C_fh)$ there.
Convolution by $f$ is bounded on $H_{J,s}$:
the radial weight has the translation bound
\[
 \|\lambda_b h\|_{H_{J,s}}
  \le2^{s/4}\max(|b|^{B_0},|b|^{-B_0})
       (1+|\log|b||^2)^{s/4}\|h\|_{H_{J,s}},
 \qquad B_0=\max(|c|,|d|),
\]
and $f$ is compactly supported and integrable.
It preserves the theta image and its closure,
so it induces a bounded operator $\overline C_f$
on the Hilbert quotient. Since $E_S$ is dense
in $H_{J,s}^{U_S}$, the preceding equality extends
to $A_f=\overline C_fq$ on the full Hilbert space.
The quotient has its isometric orthogonal
section $\sigma:Z_{J,s}^{U_S}\to H_{J,s}^{U_S}$.
Thus $\overline C_f=A_f\sigma$ is trace class.
The rank-one bounds just proved are uniform for
bounded test sets, and extension by zero to
$(Z_{J,s}^{U_S})^\perp$ preserves the trace norm.

\emph{The full trace and absolute character summation.}
All compact character projectors commute with
$C_f$, so (M27.6) descends to an orthogonal
sum of quotient sectors. Their operators are
exactly those of MT-M26, and their traces are
the inner sums in (M27.3). A trace-class operator
on an orthogonal sum has trace equal to the
sum of its sector traces, by computing in
orthonormal bases of the summands.
For clarity, absolute convergence of the full
\emph{zero} sum also follows directly from its
Hilbert trace norm. Take any finite collection
of the finite block subspaces of MT-M26 in
finitely many character sectors. Their direct
sum is finite dimensional and invariant under
$C_f$; its eigenvalues are exactly the indicated
Mellin values with their block multiplicities.
An orthonormal Schur basis for this invariant
subspace gives
\[
       \sum_{\text{these blocks}}m_{\chi,\rho}
            |\widehat f(\beta_\chi|\cdot|^\rho)|
                     \le\|C_f\|_1.
                                                               \tag{M27.12}
\]
Indeed choose diagonal phases so that the sum
of the absolute diagonal entries is a trace
against a finite-rank contraction; its absolute
value is at most the trace norm. Increasing the
finite set of blocks proves absolute convergence
of the whole spectral sum and the same bound.
Here are also the tails needed for uniform absolute
convergence, rather than only a uniform bound on the
total sum. A bounded test set lies in one space of
smooth functions with fixed compact support on an
elementary quotient. It is relatively compact there
by the derivative bounds and Arzel\`a--Ascoli.
The linear map from this Fr\'echet test space to
trace-class operators is continuous: its boundedness
on bounded sets, just proved, implies continuity
into a Banach space. Explicitly, failure of continuity
would give a sequence tending to zero in successively
more of the defining seminorms while its image norms
tend to infinity, contradicting bounded-set preservation.
The image of the test set is consequently relatively
compact in trace norm. If $P_F$ is a finite sum of
compact-character projections, then $P_F\to1$
strongly; hence
$\|T-P_FTP_F\|_1\to0$ uniformly for $T$ in this
relatively compact trace-class set. Pointwise this
follows by approximation by finite-rank operators;
uniformity follows from a finite trace-norm net and
the contraction property of the projections.
Applying (M27.12) on the complementary sectors bounds
all their absolute spectral terms by this tail norm.
For the remaining finitely many characters, the
fixed-sector estimates of MT-M26 give uniform
decay of the absolute sums outside finitely many
height blocks. The character projection of a
bounded test set is bounded in that sector's
$\mathcal E$, so those estimates apply. Combining
the two tails proves the asserted uniform absolute
convergence on bounded test sets.
This proves (M27.3).
\end{proof}

For the centered or otherwise twisted literal Masullo
quotient, put $J=a+I$. Multiplication by $|\cdot|^a$
is an isometry from (M27.1) to the Hilbert space with
base $I$ and denominator
$\overline{|\cdot|^a\Theta_K}$, and
$C_kM_a=M_aC_{|\cdot|^{-a}k}$.
Consequently, whenever $[0,1]\subset a+\operatorname{int}I$,
its full trace is
\[
       \sum_{\chi,\rho}m_{\chi,\rho}
                    \widehat k(\beta_\chi|\cdot|^{\rho-a}).
                                                               \tag{M27.13}
\]
For $a=1/2$ and $I=[-\alpha,\alpha]$, the unconditional
sufficient condition is $\alpha>1/2$.

At this centered wide-strip quotient, let
$h^*(g)=\overline{h(g^{-1})}$. Under $M$ every spectral
quasi-character in (M27.13) is unitary, and therefore
\[
       \operatorname{Tr}C_{h*h^*}
          =\sum_{\chi,\rho}m_{\chi,\rho}
               |\widehat h(\beta_\chi|\cdot|^{\rho-1/2})|^2\ge0.
                                                               \tag{M27.14}
\]
This uses $M$ as a condition; trace class alone does
not supply this positivity. Conversely, positivity
for every $h\in\mathcal D(C)$ would imply GRH for
every Hecke character of this fixed field: a single
off-line pair gives a negative fixed-character
Schwartz test by MT-M26; multiplying it by smooth
expanding norm cutoffs approximates it in that
sector's $\mathcal E$ topology by elements of
$\mathcal D(C)$, preserving a negative trace for
large cutoff. Continuity of the fixed-sector
trace-class map justifies this last step. No
positive-operator realization is inferred.

\paragraph{What remains after this transfer.}
This is the full adelic wide-strip result, not merely
a list of fixed-character trace formulas. The
varying-character summability is supplied by
Meyer's classical uniform nuclear extension, with
an explicit map into Masullo's specified Hilbert
quotient. The all-character test topology remains
the ordinary inductive union of finite invariance
levels; constants need not be uniform when that
level is allowed to vary without a boundedness
condition. The narrow centered intervals
$0<\alpha\le1/2$, even under $M$, do not embed
into the two Hilbert spaces in (M27.5).
MT-M26 treats each of their fixed-character
sectors; simultaneous trace-class control on
such a narrow interval remains a separate gate.

The original Masullo norm and closure locators
remain (9.2.2.31), p.~416; Definition~10.1.1.1,
p.~434; and (10.4.1.3), p.~456. The geometric
crosswalk to inherited cote~1 / SA23
geo-H-lidskii and geo-H-genus now includes the
actual full adelic Hilbert quotient on a wide
strip. Meyer's classical spectral realization
and uniform nuclearity retain their original
credit; this proof supplies the precise
comparison and completion transfer used here.
